\subsubsection{Barra pivotada con partícula deslizante y resorte}

Una partícula de masa $m$ se desliza sin fricción a lo largo de una barra homogénea de masa
$M$ y longitud $L$, la cual pivota libremente en torno a una bisagra fija en el punto $O$,
como se muestra en la Figura \ref{fig:c4_e01_p07_barra_particula}. La partícula está unida al
pivote $O$ mediante un resorte lineal de constante elástica $K$ y longitud natural $r_0$. La
distancia $r$ describe la posición de la partícula respecto a $O$, el ángulo $\theta$ describe
la orientación de la barra respecto a la vertical descendente, y la gravedad actúa hacia abajo.

\begin{enumerate}
    \item Demuestre que el momento de inercia de la barra respecto al eje de giro en $O$ es
          $I_{zz}=\tfrac13 ML^2$.
    \item Determine las ecuaciones diferenciales de movimiento del sistema
          ($\ddot{\theta}$ y $\ddot{r}$).
\end{enumerate}

\begin{figure}[!ht]
    \centering
    \begin{tikzpicture}[>=Latex, scale=1.4]

        % --- PARAMETROS ---
        \def\ang{35}     % Angulo de la barra respecto a la vertical
        \def\Lbar{3.6}   % Largo de la barra
        \def\rm{2.4}     % Distancia O a particula

        % --- SOPORTE FIJO (techo) ---
        \fill[pattern=north east lines] (-0.9,0) rectangle (0.4,0.18);
        \draw[thick] (-0.9,0) -- (0.4,0);

        % --- PIVOTE O ---
        \coordinate (O) at (0,0);

        % --- DIRECCION DE LA BARRA ---
        \coordinate (End) at ({\Lbar*sin(\ang)}, {-\Lbar*cos(\ang)});
        \coordinate (Pm)  at ({\rm*sin(\ang)},  {-\rm*cos(\ang)});

        % --- BARRA (cuerpo rigido) ---
        \draw[line width=5pt, line cap=round, gray!55] (O) -- (End);
        \node[below right] at (End) {$M,\,L$};

        % --- RESORTE de O a la particula ---
        \draw[thick, decoration={coil, aspect=0.5, segment length=4pt, amplitude=3.2pt, pre length=3pt, post length=3pt}, decorate] (O) -- (Pm);
        \node at ({0.55*sin(\ang)-0.35},{-0.55*cos(\ang)}) {$K$};

        % --- PARTICULA m ---
        \filldraw[fill=gris1!30, thick] (Pm) circle (3.2pt);
        \node[below left=1pt] at (Pm) {$m$};

        % --- PIVOTE ---
        \filldraw (O) circle (1.6pt) node[above left=-1pt] {$O$};

        % --- VERTICAL DE REFERENCIA Y ANGULO ---
        \draw[dashed] (O) -- (0,-2.4);
        \draw[->] (0,-1.2) arc (-90:{-90+\ang}:1.2);
        \node at ({-90+\ang/2}:1.45) {$\theta$};

        % --- COTA r (paralela a la barra) ---
        \draw[<->] ($(O)+({0.42*cos(\ang)},{0.42*sin(\ang)})$) -- ($(End)+({0.42*cos(\ang)},{0.42*sin(\ang)})$) node[midway, above right] {$r$};

        % --- GRAVEDAD ---
        \draw[->, thick] (2.3,-0.2) -- (2.3,-1.0) node[right, midway] {$\mathbf{g}$};

    \end{tikzpicture}
    \caption{Barra homogénea $M,L$ pivotada en $O$ con una partícula $m$ que desliza sobre ella, ligada a $O$ por un resorte.}
    \label{fig:c4_e01_p07_barra_particula}
\end{figure}
